% ============================================================
\setpartpreamble{\partintro
  {Eingreifen beim Speichern, Verbinden und Ausführen, und Ereignisse sauber melden.}
  {
    \item Kapitel~\ref{ch:interceptors}: Interceptors
    \item Kapitel~\ref{ch:diagnostics-and-logging}: Diagnose und Logging
  }}
\part{Laufzeit}\label{part:runtime}

% ------------------------------------------------------------
\chapter{Interceptors}\label{ch:interceptors}
\versionrisk{stable}

Interceptors sind Haken in EFs Laufzeit: vor und nach Befehlen, Verbindungen, Transaktionen,
SaveChanges und der Materialisierung. Sie können beobachten, Werte ändern, Operationen
unterdrücken oder Ergebnisse ersetzen.

\begin{datatable}{H{1.00}H{1.50}H{0.50}}
  \textbf{Interface} & \textbf{Ereignisse} & \textbf{Lebensdauer} \\
  \tablehead
  \code{ISave\allowbreak{}Changes\allowbreak{}Interceptor} & SavingChanges, SavedChanges,
  SaveChangesFailed, ThrowingConcurrencyException & Scoped \\
  \tablerule
  \code{IDb\allowbreak{}Connection\allowbreak{}Interceptor} & Öffnen, Schließen, Erzeugen
  der Verbindung & Scoped \\
  \tablerule
  \code{IDb\allowbreak{}Command\allowbreak{}Interceptor} & Erzeugen und Ausführen von
  Befehlen, DataReader & Scoped \\
  \tablerule
  \code{IDb\allowbreak{}Transaction\allowbreak{}Interceptor} & Beginnen, Commit, Rollback,
  Savepoints & Scoped \\
  \tablerule
  \code{IIdentity\allowbreak{}Resolution\allowbreak{}Interceptor} & Konflikte beim Anhängen
  bereits getrackter Entitäten; nie bei Abfragen (Kapitel~\ref{ch:materialization-from-rows-to-objects}) & Scoped \\
  \tablerule
  \code{IMaterialization\allowbreak{}Interceptor} & Erzeugen und Initialisieren von
  Entitäten (Kapitel~\ref{ch:materialization-from-rows-to-objects}) & \textbf{Singleton} \\
  \tablerule
  \code{IQuery\allowbreak{}Expression\allowbreak{}Interceptor} & LINQ-Ausdruck vor der
  Kompilierung & \textbf{Singleton} \\
\end{datatable}

Die beiden Singleton-Interceptors implementieren \code{ISingleton\allowbreak{}Interceptor}.
Sie werden einmal pro internem Provider verwendet, nicht pro Context, und müssen deshalb
zustandslos und thread-sicher sein. Im internen Provider registriert, findet EF aber nur den
Materialisierungs-Interceptor unter \code{ISingleton\allowbreak{}Interceptor}; einen
Query-Expression-Interceptor ruft es nur auf, wenn er als \code{IInterceptor} registriert ist
(Kapitel~\ref{ch:bulk-operations-and-the-limits}).

\section{Zwei Wege der Registrierung}

\begin{datatable}{H{0.96}H{1.20}H{0.84}}
  \textbf{Weg} & \textbf{Eigenschaften} & \textbf{Wann} \\
  \tablehead
  \code{options\allowbreak{}Builder.\allowbreak{}Add\allowbreak{}Interceptors(\allowbreak{}.\allowbreak{}.\allowbreak{}.\allowbreak{})}
  & Instanzen aus der Anwendung; können App-Dienste nutzen; pro Context verschieden &
  Interceptors mit App-Abhängigkeiten (Mandant, Benutzer, Uhr) \\
  \tablerule
  \code{services.\allowbreak{}Add\allowbreak{}Scoped<\allowbreak{}IInterceptor,\allowbreak{}
  X>(\allowbreak{})} in \code{Apply\allowbreak{}Services} & Im internen Provider aufgelöst
  („injected interceptors“); haben Zugriff auf EF-Dienste, nicht auf App-Dienste & Rein
  technische Interceptors ohne App-Abhängigkeit \\
\end{datatable}

Beide Wege kombiniert EF zu einer Liste. Das Beispiel nutzt den ersten, weil beide
Interceptors den Mandanten der Anwendung brauchen
(Kapitel~\ref{ch:public-api-and-integration-into}).

\section{Der SaveChanges-Interceptor}

Er setzt die \code{Tenant\allowbreak{}Id} neuer Entitäten und verhindert, dass Daten eines
anderen Mandanten geschrieben oder umgehängt werden. Das ergänzt RLS (Defense in Depth) und
liefert verständlichere Fehler als die Datenbank.
Kapitel~\ref{ch:bulk-operations-and-the-limits} zeigt, welche Wege an ihm vorbeiführen:

% source: src/Acme.EntityFrameworkCore.Tenancy/Interceptors/TenantSaveChangesInterceptor.cs
\begin{sourcecode}[Interceptors/\allowbreak{}Tenant\allowbreak{}Save\allowbreak{}Changes\allowbreak{}Interceptor.\allowbreak{}cs (\allowbreak{}1/\allowbreak{}2)]{csharp}
public class TenantSaveChangesInterceptor : SaveChangesInterceptor
{
    private readonly ITenantProvider _tenants;

    public TenantSaveChangesInterceptor(ITenantProvider tenants)
    {
        ArgumentNullException.ThrowIfNull(tenants);

        _tenants = tenants;
    }

    public override InterceptionResult<int> SavingChanges(
        DbContextEventData eventData, InterceptionResult<int> result)
    {
        if (eventData.Context is { } context) ApplyTenant(context);
        return result;
    }

    public override ValueTask<InterceptionResult<int>> SavingChangesAsync(
        DbContextEventData eventData, InterceptionResult<int> result,
        CancellationToken cancellationToken = default)
    {
        if (eventData.Context is { } context) ApplyTenant(context);
        return ValueTask.FromResult(result);
    }
\end{sourcecode}

\begin{sourcecode}[Interceptors/\allowbreak{}Tenant\allowbreak{}Save\allowbreak{}Changes\allowbreak{}Interceptor.\allowbreak{}cs (\allowbreak{}2/\allowbreak{}2,\allowbreak{} vereinfacht)]{csharp}
    internal void ApplyTenant(DbContext context)
    {
        // Welche Entitäten betroffen sind, steht in der Laufzeit-Annotation
        // (Kapitel „Fluent API auf Property-Ebene und Laufzeit-Annotations“)
        var changes = context.ChangeTracker.Entries()
            .Where(e => e.State is EntityState.Added or EntityState.Modified
                                   or EntityState.Deleted)
            .Select(e => (Entry: e, Name: e.Metadata.GetTenantPropertyName()))
            .Where(x => x.Name is not null)
            .Select(x => (x.Entry, Property: x.Entry.Property(x.Name!)))
            .ToList();
        if (changes.Count == 0) return;

        var current = _tenants.CurrentTenantId
            ?? throw new TenancyException("Kein aktiver Mandant ...");

        foreach (var (entry, property) in changes)
        {
            var currentValue = (TenantId)property.CurrentValue!;
            var originalValue = (TenantId)property.OriginalValue!;
            switch (entry.State)
            {
                case EntityState.Added when currentValue.IsEmpty:
                    property.CurrentValue = current;
                    break;
                case EntityState.Added when currentValue != current:
                    throw Mismatch(entry, currentValue, current);
                // Nur ein tatsächlicher Wertwechsel zählt, nicht die Markierung
                case EntityState.Modified
                    when property.IsModified && currentValue != originalValue:
                    throw new TenancyException("Mandanten-Property ist unveränderlich.");
                case EntityState.Modified or EntityState.Deleted:
                    if (entry.State == EntityState.Modified)
                        property.IsModified = false;    // Spalte nicht schreiben
                    // Leer = unbekannt (Shadow-Property): die Datenbank entscheidet
                    if (!originalValue.IsEmpty && originalValue != current)
                        throw Mismatch(entry, originalValue, current);
                    break;
            }
        }
    }
}
\end{sourcecode}

Wer nur mit asynchronem Code arbeitet, überschreibt leicht nur eine Variante:

\begin{sourcecode}[Interceptor (\allowbreak{}FALSCH)]{csharp}
public class TenantSaveChangesInterceptor : SaveChangesInterceptor
{
    public override ValueTask<InterceptionResult<int>> SavingChangesAsync(
        DbContextEventData eventData, InterceptionResult<int> result,
        CancellationToken cancellationToken = default) { /* Mandant setzen */ }

    // SavingChanges fehlt
}
\end{sourcecode}

\begin{failure}[Was passiert, wenn nur die asynchrone Variante implementiert ist?]
  \textbf{Beobachtbar.} Mit \code{Save\allowbreak{}Changes\allowbreak{}Async} funktioniert
  alles. Ein synchrones \code{Save\allowbreak{}Changes}, etwa in einem Konsolenwerkzeug oder
  einem älteren Codepfad, speichert ohne Mandant; RLS lehnt die Zeile mit
  \code{WITH CHECK} ab, ein Plugin ohne Datenbankschutz schreibt sie still falsch.

  \textbf{Ursache.} EF ruft je nach Aufruf genau eine Variante auf. Die Basisklasse
  \code{Save\allowbreak{}Changes\allowbreak{}Interceptor} implementiert die andere als
  leere Methode, also ohne Fehler. Dasselbe gilt für Connection- und Command-Interceptors.

  \textbf{Lösung.} Immer beide Varianten überschreiben und dieselbe Logik aufrufen, wie
  \code{Apply\allowbreak{}Tenant} im Beispiel.

  \textbf{Test:}
  \code{Failure\allowbreak{}Case\allowbreak{}Tests.\allowbreak{}Interceptor\allowbreak{}\_with\allowbreak{}\_only\allowbreak{}\_the\allowbreak{}\_async\allowbreak{}\_variant\allowbreak{}\_…}.
\end{failure}

\section{Markierung, Wertwechsel und Datenbankschutz}\label{sec:tracking-flag-versus-value}

Drei Dinge sehen beim Speichern ähnlich aus und bedeuten Verschiedenes:

\begin{itemize}
  \item \textbf{Die Tracking-Markierung.} \code{Db\allowbreak{}Context.\allowbreak{}Update} und
    \code{Attach} einer nicht getrackten Entität setzen \code{Is\allowbreak{}Modified} für
    \emph{jede} Property, auch für die Mandanten-Property. Original- und aktueller Wert sind
    dabei gleich, nämlich der Wert, den die Entität mitbringt. Bis 1.0 lehnte der Interceptor
    schon diese Markierung als Mandantenwechsel ab, und ein gewöhnliches Update aus einer API
    scheiterte.
  \item \textbf{Der tatsächliche Wertwechsel.} Nur wenn sich der Wert gegenüber dem
    Originalwert ändert, wird eine Zeile umgehängt. Das lehnt der Interceptor ab, auch wenn
    der neue Wert der aktive Mandant ist. Bei einem unveränderten Wert nimmt er die Markierung
    zurück, sodass das \code{UPDATE} die Mandantenspalte nicht schreibt.
  \item \textbf{Der Datenbankzustand.} Bei einer getrackten Entität stammt der Originalwert
    aus der Datenbank. Bei einer abgekoppelten Entität ist er nur eine Behauptung des
    Aufrufers. Eine manipulierte Entität mit dem eigenen Mandanten und der ID einer fremden Zeile
    besteht deshalb jede Prüfung im Interceptor. Erst die Policy entscheidet: Das
    \code{UPDATE} sieht die fremde Zeile nicht, ändert null Zeilen, und EF meldet eine
    \code{Db\allowbreak{}Update\allowbreak{}Concurrency\allowbreak{}Exception}. Ein leerer
    Originalwert (eine Shadow-Property, die die Entität nicht tragen kann) gilt als unbekannt,
    nicht als fremd; auch hier entscheidet die Datenbank.
\end{itemize}

Der Interceptor verspricht also nur, was er ohne Datenbankzugriff wissen kann: keinen
Mandantenwechsel im Change Tracker, keine bekannt fremden Werte, keine Änderung ohne aktiven
Mandanten. Ob eine Zeile wirklich dem Mandanten gehört, prüft allein RLS.
\code{Save\allowbreak{}Changes\allowbreak{}Interceptor\allowbreak{}Tests} deckt die Fälle
ohne Datenbank ab, \code{Detached\allowbreak{}Update\allowbreak{}Tests} gegen PostgreSQL mit
eingeschränkter Rolle: \code{Update} mit unverändertem Mandanten synchron und asynchron, die
manipulierte fremde Zeilen-ID und das Umhängen einer getrackten Zeile.

Die Logik liegt in einer \code{internal}-Methode \code{Apply\allowbreak{}Tenant}. Tests
können sie über \code{Internals\allowbreak{}Visible\allowbreak{}To} direkt aufrufen, ohne
eine Datenbank zu brauchen (Kapitel~\ref{ch:test-strategy}).

\section{Der Connection-Interceptor}

Er setzt bei jedem Öffnen der Verbindung das Session-Setting, das die Policies und
\code{Current\allowbreak{}Tenant(\allowbreak{})} auswerten. Der Ablauf (Öffnen abfangen,
sync und async, Mandant holen) ist für jede Datenbank gleich und liegt deshalb in einer
abstrakten Basis im relationalen Paket (Kapitel~\ref{ch:what-belongs-in-the-core}). Das
PostgreSQL-Paket liefert nur den Befehl. Den Setting-Namen holt es über den Context aus den
Event-Daten
(\code{context.\allowbreak{}Get\allowbreak{}Service<\allowbreak{}T>(\allowbreak{})}) aus dem
\emph{internen} Provider, den Mandanten aus dem \emph{App}-Container:

\begin{sourcecode}[Interceptors/\allowbreak{}Tenant\allowbreak{}Session\allowbreak{}Interceptor.\allowbreak{}cs (\allowbreak{}relational),\allowbreak{} Runtime/\allowbreak{}Tenant\allowbreak{}Connection\allowbreak{}Interceptor.\allowbreak{}cs (\allowbreak{}vereinfacht)]{csharp}
// Relationales Paket
public abstract class TenantSessionInterceptor : DbConnectionInterceptor
{
    protected TenantSessionInterceptor(ITenantProvider tenants)
    {
        ArgumentNullException.ThrowIfNull(tenants);

        Tenants = tenants;
    }

    protected ITenantProvider Tenants { get; }

    public override async Task ConnectionOpenedAsync(
        DbConnection connection, ConnectionEndEventData eventData,
        CancellationToken cancellationToken = default)
    {
        await using var command =
            CreateSessionCommand(connection, eventData.Context, Tenants.CurrentTenantId);
        await command.ExecuteNonQueryAsync(cancellationToken);
    }

    // ConnectionOpened (synchron) analog

    protected abstract DbCommand CreateSessionCommand(
        DbConnection connection, DbContext? context, TenantId? tenant);
}

// PostgreSQL-Paket
public class TenantConnectionInterceptor : TenantSessionInterceptor
{
    public TenantConnectionInterceptor(ITenantProvider tenants)
        : base(tenants)
    {
    }

    protected override DbCommand CreateSessionCommand(
        DbConnection connection, DbContext? context, TenantId? tenant)
    {
        var settingName = context?.GetService<ITenancySingletonOptions>()
                              .Settings.SettingName
                          ?? TenancySettings.Default.SettingName;

        var command = connection.CreateCommand();
        command.CommandText = "SELECT set_config(@name, @value, false)";
        AddParameter(command, "name", settingName);
        // Immer setzen, auch leer: kein alter Wert bleibt stehen
        AddParameter(command, "value", tenant?.ToString() ?? string.Empty);
        return command;
    }
}
\end{sourcecode}

\begin{note}[Leerer String statt NULL]
  Npgsql setzt den Session-Zustand zurück, wenn eine Verbindung in den Pool zurückkehrt.
  Danach liefert
  \code{current\_\allowbreak{}setting(\allowbreak{}'app.\allowbreak{}tenant\_\allowbreak{}id',\allowbreak{}
  true)} einen leeren String statt \code{NULL}. Ein \code{''::uuid} würde jede Abfrage
  abbrechen. Policy und Translator verwenden deshalb
  \code{nullif(\allowbreak{}.\allowbreak{}.\allowbreak{}.\allowbreak{},\allowbreak{} '')}:
  Ohne Mandant ist der Wert \code{NULL}, der Vergleich nie wahr, und es sind keine Zeilen
  sichtbar. Das ist die sichere Richtung.
\end{note}

EF öffnet und schließt die Verbindung standardmäßig für jede Abfrage. Jedes Öffnen kostet
hier einen zusätzlichen Roundtrip. Bei vielen kleinen Abfragen lohnt es sich, die Verbindung
pro Anfrage einmal explizit zu öffnen (Kapitel~\ref{ch:plugin-performance}).

\section{Wo der Connection-Interceptor nicht läuft}

Der Interceptor läuft, wenn EF eine Verbindung öffnet, und nur dann. Zwei Situationen
umgehen ihn, und beide entstehen leicht aus Versehen:

\begin{datatable}{H{1.00}H{1.06}H{0.94}}
  \textbf{Situation} & \textbf{Folge} & \textbf{Abhilfe} \\
  \tablehead
  Eine bereits geöffnete Verbindung wird übergeben
  (\code{Use\allowbreak{}Npgsql(\allowbreak{}connection)},
  \code{Set\allowbreak{}Db\allowbreak{}Connection}) & EF öffnet sie nie, also läuft der
  Interceptor nie. Die Session behält ihren bisherigen Wert: keinen oder den eines früheren
  Mandanten & Verbindungen geschlossen übergeben und EF öffnen lassen; sonst
  \code{Database.\allowbreak{}Apply\allowbreak{}Tenant\allowbreak{}Session(\allowbreak{})}
  aufrufen \\
  \tablerule
  Die Verbindung bleibt offen (\code{Open\allowbreak{}Connection}, eine offene Transaktion),
  während der Mandant wechselt & Der neue Mandant wirkt erst, wenn die Verbindung erneut
  geöffnet wird & Ein Context, ein Mandant; ist das nicht möglich, nach dem Wechsel
  \code{Database.\allowbreak{}Apply\allowbreak{}Tenant\allowbreak{}Session(\allowbreak{})}
  aufrufen \\
\end{datatable}

\begin{sourcecode}[tests/\allowbreak{}.\allowbreak{}.\allowbreak{}.\allowbreak{}Integration\allowbreak{}Tests/\allowbreak{}Connection\allowbreak{}Lifetime\allowbreak{}Tests.\allowbreak{}cs (\allowbreak{}gekürzt, vereinfacht)]{csharp}
// Bereits offen, mit einem Session-Wert aus früherer Arbeit (Mandant B)
await using var connection = new NpgsqlConnection(appConnectionString);
await connection.OpenAsync();
// ... SELECT set_config('app.tenant_id', <tenant B>, false) ...

// Gedacht für Mandant A, aber EF öffnet die Verbindung nicht: kein Interceptor
await using var db = CreateAppContext(connection, tenantA);
var orders = await db.Orders.ToListAsync();       // die Zeilen von Mandant B!
\end{sourcecode}

Für diese Fälle bietet das Plugin das explizite Gegenstück zum Interceptor.
\code{Database.\allowbreak{}Apply\allowbreak{}Tenant\allowbreak{}Session(\allowbreak{})}
schreibt den aktuellen Mandanten in die Session, wenn die Verbindung offen ist, und tut
nichts, wenn sie geschlossen ist, denn dann öffnet EF sie und der Interceptor übernimmt. Die
Methode verwendet den Befehl des Interceptors selbst und funktioniert deshalb mit jedem
Provider-Paket:

\begin{sourcecode}[Verwendung]{csharp}
await using var db = CreateAppContext(connection, tenantA);   // Verbindung bereits offen
await db.Database.ApplyTenantSessionAsync();                  // Session jetzt: Mandant A
var orders = await db.Orders.ToListAsync();                   // die Zeilen von Mandant A
\end{sourcecode}

% source: src/Acme.EntityFrameworkCore.Tenancy.Relational/Interceptors/TenantSessionInterceptor.cs
\begin{sourcecode}[Interceptors/\allowbreak{}Tenant\allowbreak{}Session\allowbreak{}Interceptor.\allowbreak{}cs (\allowbreak{}Auszug)]{csharp}
public virtual async Task<bool> ApplyToOpenConnectionAsync(
    DbContext context, CancellationToken cancellationToken = default)
{
    var connection = context.Database.GetDbConnection();
    if (connection.State != ConnectionState.Open)
        return false;                     // EF öffnet sie später; der Interceptor übernimmt

    await using var command =
        CreateSessionCommand(connection, context, Tenants.CurrentTenantId);
    // In einer expliziten Transaktion muss der Befehl ihr zugeordnet sein
    command.Transaction = context.Database.CurrentTransaction?.GetDbTransaction();
    await command.ExecuteNonQueryAsync(cancellationToken);
    return true;
}
\end{sourcecode}

Die Erweiterungsmethode \code{Apply\allowbreak{}Tenant\allowbreak{}Session} im relationalen Paket
sucht den Interceptor unter den Interceptoren des Contexts und bricht mit einer klaren
Meldung ab, wenn keiner registriert ist. Zwei Details sind wichtig: Die Methode braucht den
Context, nicht nur die Verbindung, weil der Provider seine Einstellungen aus dem internen
Service Provider liest; deshalb erhält \code{Create\allowbreak{}Session\allowbreak{}Command}
den \code{Db\allowbreak{}Context} statt der Event-Daten. Und der Befehl muss sich einer
laufenden Transaktion anschließen: PostgreSQL ist hier nachsichtig, SQL Server lehnt einen
Befehl ab, der ihr nicht zugeordnet ist. Die SQL-Server-Tests prüfen genau diesen Fall; ohne
die Zeile schlagen sie fehl.

In \code{Connection\allowbreak{}Opened} fehlt diese Zuordnung mit Absicht. Eine lokale
Transaktion setzt eine offene Verbindung voraus, und EF öffnet die Verbindung, bevor es
\code{Begin\allowbreak{}Transaction} aufruft; beim Öffnen gibt es also noch keine Transaktion.
Eine umgebende \code{Transaction\allowbreak{}Scope} wiederum übernehmen Npgsql und SqlClient
beim Öffnen selbst, und jeder Befehl auf der Verbindung läuft darin, ohne dass ihm eine
\code{Db\allowbreak{}Transaction} zugewiesen wird. \code{Transaction\allowbreak{}Context\allowbreak{}Tests}
und die SQL-Server-Tests prüfen alle drei Wege: die Verbindung, die
\code{Begin\allowbreak{}Transaction\allowbreak{}Async} öffnet,
\code{Apply\allowbreak{}Tenant\allowbreak{}Session} in einer expliziten Transaktion und
Lesen und Schreiben unter einer \code{Transaction\allowbreak{}Scope}.

\begin{important}[Ein Context, ein Mandant]
  Die Isolation folgt der Verbindung, nicht dem Context. Ein Context, der für mehrere
  Mandanten verwendet wird, oder eine Verbindung, die eine Anfrage überdauert, trägt den
  Mandanten des Moments, in dem sie geöffnet wurde.
  \code{Connection\allowbreak{}Lifetime\allowbreak{}Tests} zeigt beide Fälle gegen eine
  echte Datenbank, zusammen mit den beiden Abhilfen: einer geschlossenen Verbindung, die EF
  selbst öffnet, und \code{Apply\allowbreak{}Tenant\allowbreak{}Session(\allowbreak{})}. Ein
  Plugin, dessen Hooks nicht jeden Weg abdecken, sollte eine solche explizite API anbieten,
  statt die Lücke nur zu dokumentieren.
\end{important}

Erkennen lässt sich die Lücke zusätzlich beim Laden: Mit
\code{Verify\allowbreak{}Materialized\allowbreak{}Tenants(\allowbreak{})} prüft das Plugin jede
geladene Entität gegen den Mandanten des Contexts und wirft, statt fremde Zeilen
zurückzugeben (Kapitel~\ref{ch:materialization-from-rows-to-objects}).

\section{Dapper und rohes ADO.NET}

Reports und Massenexporte laufen oft über Dapper oder rohes ADO.NET auf der Verbindung des
Contexts (\code{Database.\allowbreak{}Get\allowbreak{}Db\allowbreak{}Connection(\allowbreak{})}).
Row Level Security gilt auch für diese Befehle, weil sie in der Datenbank liegt und nicht in EF.
Es kommt nur darauf an, ob der Session-Wert gesetzt ist, und das hängt davon ab, wer die
Verbindung öffnet:

\begin{datatable}{H{0.95}H{1.05}}
  \textbf{Wer die Verbindung öffnet} & \textbf{Ergebnis} \\
  \tablehead
  EF, etwa \code{Open\allowbreak{}Connection\allowbreak{}Async(\allowbreak{})} oder eine offene
  Transaktion & Der Interceptor hat die Session gesetzt; Dapper sieht die Zeilen des aktuellen
  Mandanten \\
  \tablerule
  Dapper oder ADO.NET selbst, auch nachdem EF die Verbindung nach einer Abfrage wieder
  geschlossen hat & Kein Session-Wert: Die Policies treffen keine Zeile (fail closed), bis
  \code{Apply\allowbreak{}Tenant\allowbreak{}Session\allowbreak{}Async(\allowbreak{})} ihn setzt \\
  \tablerule
  Mit Transaction Pooling (\code{Use\allowbreak{}Transaction\allowbreak{}Pooling}) & Der Mandant
  existiert nur in der Transaktion von EF; Dapper muss darin laufen \\
\end{datatable}

Die zweite Zeile ist die Falle: EF öffnet und schließt die Verbindung für jede Abfrage, eine
Verbindung, die eben noch den Mandanten trug, ist danach also wieder geschlossen. Dapper öffnet
eine geschlossene Verbindung selbst, und der Interceptor läuft nicht. Das sichere Muster öffnet
die Verbindung für den ganzen Block über EF:

\begin{sourcecode}[Aufruf (\allowbreak{}Dapper)]{csharp}
await db.Database.OpenConnectionAsync();                // EF öffnet: Session gesetzt
var connection = db.Database.GetDbConnection();
var orders = await connection.QueryAsync<OrderDto>(
    "SELECT \"Id\", \"Total\" FROM \"Orders\"");             // RLS: nur der Mandant
await db.Database.CloseConnectionAsync();

// PgBouncer im Transaktionsmodus: Mandant nur in der Transaktion von EF
await using var tx = await db.Database.BeginTransactionAsync();
var rows = await connection.QueryAsync<OrderDto>(sql, transaction: tx.GetDbTransaction());
\end{sourcecode}

% source: tests/Acme.EntityFrameworkCore.Tenancy.IntegrationTests/RawAdoNetTests.cs
\begin{sourcecode}[tests/\allowbreak{}.\allowbreak{}.\allowbreak{}.\allowbreak{}/\allowbreak{}Raw\allowbreak{}Ado\allowbreak{}Net\allowbreak{}Tests.\allowbreak{}cs (\allowbreak{}Auszug)]{csharp}
var tenant = await CreateTenantWithOrdersAsync(2);
await using var db = _fixture.CreateAppContext(
    new MutableTenantProvider { CurrentTenantId = tenant });
await db.Orders.CountAsync();                 // EF hat geöffnet und wieder geschlossen

// Direkt geöffnet, wie Dapper es bei geschlossener Verbindung tut: ohne EF
var connection = db.Database.GetDbConnection();
await connection.OpenAsync();
Assert.Equal(0, await CountOrdersAsync(connection));      // fail closed

await db.Database.ApplyTenantSessionAsync();
Assert.Equal(2, await CountOrdersAsync(connection));
\end{sourcecode}

\code{Raw\allowbreak{}Ado\allowbreak{}Net\allowbreak{}Tests} prüft alle drei Zeilen gegen
PostgreSQL mit einem einfachen Befehl, mehr schickt Dapper auch nicht. Das Ergebnis ist nie ein
Leck: Ohne Session-Wert sieht die Abfrage nichts. Es ist ein falsches leeres Ergebnis, und
deshalb gehört das Muster in die Richtlinien der Anwendung.

% ------------------------------------------------------------
\chapter{Diagnose und Logging}\label{ch:diagnostics-and-logging}
\versionrisk{stable}

Ein professionelles Plugin meldet Ereignisse über EFs Logging-System. Dann landen die
Meldungen in denselben Kanälen, die Nutzer ohnehin konfigurieren (\code{ILogger},
\code{Log\allowbreak{}To}, \code{Diagnostic\allowbreak{}Source}), und Warnungen lassen sich
per \code{Configure\allowbreak{}Warnings} ignorieren, protokollieren oder zu Fehlern machen.

\section{Event-IDs}

% source: src/Acme.EntityFrameworkCore.Tenancy/Diagnostics/TenancyEventId.cs
\begin{sourcecode}[Diagnostics/\allowbreak{}Tenancy\allowbreak{}Event\allowbreak{}Id.\allowbreak{}cs (\allowbreak{}gekürzt)]{csharp}
public static class TenancyEventId
{
    public const int TenancyBaseId = 1_150_000;     // eigener, dokumentierter Bereich

    private enum Id
    {
        TenantIsolationApplied = TenancyBaseId,
        TenantIsolationOptedOut,
        PoliciesWithoutRowLevelSecurity,
    }

    private static readonly string ModelPrefix = "Acme.EntityFrameworkCore.Tenancy.Model.";

    public static readonly EventId TenantIsolationApplied =
        Make(Id.TenantIsolationApplied);
    public static readonly EventId TenantIsolationOptedOut =
        Make(Id.TenantIsolationOptedOut);
    // ...

    private static EventId Make(Id id) => new((int)id, ModelPrefix + id);
}
\end{sourcecode}

EF und Provider belegen niedrige Bereiche. Ein Plugin wählt einen eigenen, hohen Bereich und
dokumentiert ihn. Die Namen folgen EFs Schema „Assembly.Kategorie.Ereignis“, damit Filter
nach Präfix funktionieren.

\section{Event-Definitionen}

Eine \code{Event\allowbreak{}Definition<\allowbreak{}T>} bündelt Event-ID, Log-Level,
Meldungsvorlage und das Warnverhalten aus den Options. EF cached seine Definitionen pro
internem Provider; das Plugin macht es genauso mit einer eigenen Klasse, die als Singleton
registriert ist (Kapitel~\ref{ch:registering-services}):

% source: src/Acme.EntityFrameworkCore.Tenancy/Diagnostics/TenancyLoggingDefinitions.cs (strings)
\begin{sourcecode}[Diagnostics/\allowbreak{}Tenancy\allowbreak{}Logging\allowbreak{}Definitions.\allowbreak{}cs (\allowbreak{}Auszug)]{csharp}
public class TenancyLoggingDefinitions
{
    private EventDefinition<string>? _tenantIsolationOptedOut;

    public virtual EventDefinition<string> LogTenantIsolationOptedOut(
        IDiagnosticsLogger logger)
        => LazyInitializer.EnsureInitialized(ref _tenantIsolationOptedOut, () =>
            Create(logger, TenancyEventId.TenantIsolationOptedOut, LogLevel.Warning,
                nameof(TenancyEventId.TenantIsolationOptedOut),
                "Entität '{entityType}' implementiert ITenantOwned, ist aber von der "
                + "Mandantenisolation ausgenommen."));

    private static EventDefinition<string> Create(
        IDiagnosticsLogger logger, EventId eventId, LogLevel level,
        string name, string message)
        => new(
            logger.Options,                       // enthält das Warnverhalten
            eventId,
            level,
            "TenancyEventId." + name,             // steht in Warning-as-Error-Meldungen
            l => LoggerMessage.Define<string>(l, eventId, message));
}
\end{sourcecode}

\section{Die Logger-Extension}

Die eigentliche Log-Methode folgt exakt dem Muster aus EFs
\code{Core\allowbreak{}Logger\allowbreak{}Extensions}: erst klassisches Logging, dann die
Ereignisdaten für DiagnosticSource und \code{Log\allowbreak{}To}:

\begin{sourcecode}[Diagnostics/\allowbreak{}Tenancy\allowbreak{}Logger\allowbreak{}Extensions.\allowbreak{}cs (\allowbreak{}Auszug, vereinfacht)]{csharp}
private static void LogEntityType(
    IDiagnosticsLogger<DbLoggerCategory.Model> diagnostics,
    EventDefinition<string> definition,
    IReadOnlyEntityType entityType)
{
    // 1. ILogger; wirft bei ConfigureWarnings(w => w.Throw(...))
    if (diagnostics.ShouldLog(definition))
        definition.Log(diagnostics, entityType.DisplayName());

    // 2. DiagnosticSource und LogTo (simple logging)
    if (diagnostics.NeedsEventData(definition, out var diagnosticSource, out var simpleLog))
    {
        var eventData = new EntityTypeEventData(definition, GenerateMessage, entityType);
        diagnostics.DispatchEventData(definition, eventData, diagnosticSource, simpleLog);
    }
}
\end{sourcecode}

Nutzer können die Warnung damit in einen Fehler verwandeln:

\begin{sourcecode}[Anwendung]{csharp}
options.UseNpgsql(cs)
       .UseAcmeTenancy()
       .ConfigureWarnings(w => w.Throw(TenancyEventId.TenantIsolationOptedOut));

// -> InvalidOperationException beim Modellbau:
//    "An error was generated for warning '...TenantIsolationOptedOut': Entität
//     'AuditEntry' ... This exception can be suppressed or logged by passing event ID
//     'TenancyEventId.TenantIsolationOptedOut' to the 'ConfigureWarnings' method ..."
\end{sourcecode}

\begin{tip}
  Sicherheitsrelevante Abweichungen gehören als \textbf{Warnung} ins Log und nicht als
  Debug-Meldung. Teams mit strengen Anforderungen eskalieren sie per
  \code{Configure\allowbreak{}Warnings} zu Fehlern; andere sehen sie zumindest beim Start.
  Das Plugin muss diese Entscheidung nicht für alle treffen.
\end{tip}

\section{OpenTelemetry: der Mandant im Trace}

Verteiltes Tracing mit OpenTelemetry baut auf
\code{System.\allowbreak{}Diagnostics.\allowbreak{}Activity} auf. EF Core startet selbst keine
Activities. Die Spans der Datenbankbefehle kommen entweder von der OpenTelemetry-Instrumentierung
für EF Core, die auf die \code{Diagnostic\allowbreak{}Source}-Ereignisse von EF hört, oder vom
Datenbanktreiber: Npgsql und SqlClient haben eigene Activity-Quellen. Beiden fehlt meist der
Mandant, die erste Frage, wenn eine Abfrage im Betrieb langsam ist oder scheitert. Mit
\code{Enrich\allowbreak{}Open\allowbreak{}Telemetry(\allowbreak{})} ergänzt das Plugin ihn als
Tag \code{tenant.\allowbreak{}id} (Standard: aus):

\begin{sourcecode}[Aufruf]{csharp}
options.UseNpgsql(cs).UseAcmeTenancy(sp, t => t.EnrichOpenTelemetry());
// SQL Server: UseSqlServer(cs, sql => sql.UseAcmeTenancy(sp).EnrichOpenTelemetry())
\end{sourcecode}

Ein Command-Interceptor taggt vor jedem Befehl \code{Activity.\allowbreak{}Current}. Er ist im
internen Service Provider registriert, eine Instanz für alle Contexts, und liest den Mandanten
aus dem Context des Befehls:

% source: src/Acme.EntityFrameworkCore.Tenancy.Relational/Interceptors/TenantActivityEnrichmentInterceptor.cs
\begin{sourcecode}[Interceptors/\allowbreak{}Tenant\allowbreak{}Activity\allowbreak{}Enrichment\allowbreak{}Interceptor.\allowbreak{}cs (\allowbreak{}Auszug)]{csharp}
protected virtual void Enrich(CommandEventData eventData)
{
    // Billige Prüfungen zuerst: Ohne aufgezeichnete Activity gibt es nichts zu tun
    if (Activity.Current is not { IsAllDataRequested: true } activity
        || eventData.Context is not { } context)
        return;

    // Der Tenant-Provider dieses Contexts reist mit seinem SaveChanges-Interceptor
    var tenant = context.GetService<IDbContextOptions>()
        .FindExtension<CoreOptionsExtension>()?.Interceptors?
        .OfType<TenantSaveChangesInterceptor>().FirstOrDefault()?.Tenants.CurrentTenantId;

    // Ohne Mandant kein Tag: Ein leerer oder Platzhalterwert wäre irreführend
    if (tenant is { } id)
        activity.SetTag(TagName, id.Value.ToString("D"));
}
\end{sourcecode}

Welche Activity das Tag erhält, hängt von der Instrumentierung ab:

\begin{datatable}{H{0.85}H{1.15}}
  \textbf{Instrumentierung} & \textbf{Getaggte Activity} \\
  \tablehead
  OpenTelemetry-Instrumentierung für EF Core & Der Span des Befehls selbst: EF~10 löst sein
  \code{Command\allowbreak{}Executing}-Ereignis, das den Span startet, vor den Interceptors
  aus \\
  \tablerule
  Nur der Treiber (Npgsql, SqlClient) & Der umgebende Span, meist der HTTP-Request: Der Treiber
  startet seinen Span erst in \code{Execute\allowbreak{}Reader}, nach dem Interceptor \\
  \tablerule
  Kein Listener & Keine; der Interceptor kehrt nach einer Prüfung zurück, ohne Allokationen \\
\end{datatable}

\code{Activity\allowbreak{}Enrichment\allowbreak{}Tests} prüft das Tag gegen PostgreSQL mit einem
\code{Activity\allowbreak{}Listener}: gesetzt, wenn eingeschaltet, fehlend, wenn
ausgeschaltet, und fehlend ohne Mandant.

\begin{note}[Mandanten-IDs in der Telemetrie]
  Eine Mandanten-ID ist ein Bezeichner und kein Geheimnis, aber Traces verlassen oft die
  Anwendung: zu einem Monitoring-Dienst, manchmal zu einem anderen Unternehmen. Ob
  Mandanten-IDs dorthin dürfen, entscheidet der Betreiber nach Datenschutzgesichtspunkten,
  deshalb ist die Option standardmäßig aus.
\end{note}
